\subsection{時間、信号、変化率}

制御器が最初に受け取るのは、連続した数式ではなく、一定間隔で並んだ測定記録である。水槽では水位 $h$ [cm]、室温制御では温度 $T$ [$^\circ$C]、移動ロボットでは基準線からの横ずれ $e$ [m] が、時刻とともに表へ記録される。サンプリング間隔を $\Delta t$ [s] とし、$k=0,1,2,\ldots$ 番目の記録時刻を
\begin{equation}
  t_k=k\Delta t
\end{equation}
と置く。たとえば $\Delta t=2\,\mathrm{s}$ なら、$t_0=0\,\mathrm{s}$、$t_1=2\,\mathrm{s}$、$t_2=4\,\mathrm{s}$ である。水位、温度、横ずれの測定値は
\begin{equation}
  h_k=h(t_k),\qquad T_k=T(t_k),\qquad e_k=e(t_k)
\end{equation}
のように添字つきの列として扱える。目標水位 $r_k$ [cm]、ポンプ指令 $u_{c,k}$ [\%]、実際の流入量 $u_k$ [L/s]、排水や追加流入を表す外乱 $d_k$ [L/s]、水位計の測定ばらつき $n_k$ [cm] も同じ時刻列に並ぶ。

次の入力を選ぶには、値そのものだけでなく、値がどの速さで変わっているかが必要になる。連続する二つの水位記録が
\begin{equation}
  h_k=18.0\,\mathrm{cm},\qquad h_{k+1}=18.6\,\mathrm{cm},\qquad
  \Delta t=2.0\,\mathrm{s}
\end{equation}
であったとする。この区間で増えた水位は
\begin{equation}
  \Delta h_k=h_{k+1}-h_k
  =18.6\,\mathrm{cm}-18.0\,\mathrm{cm}
  =0.6\,\mathrm{cm}
\end{equation}
である。これを経過時間で割ると、区間 $[t_k,t_{k+1}]$ における平均変化率
\begin{equation}
  \frac{\Delta h_k}{\Delta t}
  =\frac{h_{k+1}-h_k}{\Delta t}
  =\frac{0.6\,\mathrm{cm}}{2.0\,\mathrm{s}}
  =0.30\,\mathrm{cm/s}
\end{equation}
を得る。温度で同じ計算をすれば単位は $^\circ$C/s になり、ロボットの横ずれで同じ計算をすれば単位は m/s になる。分子の単位を時間の単位で割るため、変化率の単位は「その量の単位毎秒」になる。

この段階で、時刻ごとに値をもつ量を一つの対象として扱う必要が生じる。水位、温度、横ずれ、目標値、入力指令、外乱、測定ばらつきは役割が異なるが、いずれも時刻を指定すると値が定まる。
\begin{definition}[信号]
各時刻 $t$ に対して数、ベクトル、または行列を対応させるものを信号という。水位 $h(t)$、温度 $T(t)$、角度 $\theta(t)$、電圧 $v(t)$、位置 $p(t)\in\R^3$ はすべて信号である。
\end{definition}

平均変化率は測定間隔に依存する。区間を短くすると、より短い時間幅での変化を評価することになる。一般の信号 $x$ について、$t_k$ から $t_{k+1}$ までの平均変化率は
\begin{equation}
  \frac{x_{k+1}-x_k}{\Delta t}
\end{equation}
である。ここで $x_k=x(t_k)$、$x_{k+1}=x(t_k+\Delta t)$ だから、同じ式は
\begin{equation}
  \frac{x(t_k+\Delta t)-x(t_k)}{\Delta t}
\end{equation}
と書ける。サンプリング間隔を十分小さくした極限として、時刻 $t$ そのものにおける瞬間的な変化率を表す記法が微分である。

角度 $\theta$ [rad] を例にとる。時刻 $t$ から $t+\Delta t$ までの平均角速度は
\begin{equation}
  \frac{\theta(t+\Delta t)-\theta(t)}{\Delta t}
\end{equation}
であり、単位は rad/s である。$\Delta t$ を 0 に近づけた極限が存在するとき、角速度 $\omega$ を
\begin{equation}
  \omega(t)
  =\lim_{\Delta t\to0}
  \frac{\theta(t+\Delta t)-\theta(t)}{\Delta t}
  =\frac{\dd \theta(t)}{\dd t}
  \label{eq:angular-speed-definition}
\end{equation}
と書く。\weq{angular-speed-definition} は、有限差分で計算していた平均角速度を、連続時間の瞬間角速度へ移した表記である。

逆に、変化率を時間ごとに足し合わせると、量そのものの変化が得られる。区間 $[t_k,t_{k+1}]$ で水位の変化率を $v_k$ [cm/s] と近似すれば、
\begin{equation}
  h_{k+1}-h_k=v_k\Delta t
\end{equation}
である。したがって
\begin{align}
  h_1&=h_0+v_0\Delta t,\\
  h_2&=h_1+v_1\Delta t
      =h_0+v_0\Delta t+v_1\Delta t,\\
  h_3&=h_2+v_2\Delta t
      =h_0+v_0\Delta t+v_1\Delta t+v_2\Delta t
\end{align}
となる。$N$ 個の区間を通した水位は
\begin{equation}
  h_N=h_0+\sum_{k=0}^{N-1}v_k\Delta t
\end{equation}
で表される。各項 $v_k\Delta t$ の単位は $(\mathrm{cm/s})\cdot\mathrm{s}=\mathrm{cm}$ であり、総和全体も水位の単位 cm をもつ。

角速度でも同じ加算が成り立つ。$\omega_k$ [rad/s] を区間 $[t_k,t_{k+1}]$ の角速度の代表値とすれば
\begin{align}
  \theta_{k+1}-\theta_k&=\omega_k\Delta t,\\
  \theta_N&=\theta_0+\sum_{k=0}^{N-1}\omega_k\Delta t
\end{align}
である。ここで $t_N=t$ とし、$\Delta t$ を小さくしていくと、有限和は連続時間の積分に移る。初期角度 $\theta(0)$ と角速度 $\omega(t)$ が分かっているとき、
\begin{equation}
  \theta(t)
  =\theta(0)+
  \lim_{\Delta t\to0}\sum_{k=0}^{N-1}\omega(t_k)\Delta t
  =\theta(0)+\int_0^t \omega(\tau)\,\dd \tau
  \label{eq:integral-angle}
\end{equation}
である。\weq{integral-angle} は、サンプリングされた角速度を有限和で積み上げる計算を、連続時間の極限として書いたものである。微分は短い区間の差分から瞬間の変化率を取り出し、積分は短い区間の増分を足し戻す。以後のモデル式は、この二つを使って、測定記録から時間変化を表す方程式へ移る。

\subsubsection{指数関数解の導出}

温度差、速度誤差、電気回路の電流のように「現在大きいほど速く減る」量は多い。このとき最小のモデルは
\begin{equation}
  \dot{x}(t)=-a x(t),\qquad a>0
  \label{eq:basic-decay-ode}
\end{equation}
である。ここで点は時間微分を表す。

\begin{theorem}[一次の同次線形微分方程式の解]
\weq{basic-decay-ode} の解は
\begin{equation}
  x(t)=x(0)e^{-at}
  \label{eq:basic-decay-solution}
\end{equation}
である。
\end{theorem}

\begin{derivation}
$x(t)\ne0$ の区間で両辺を $x(t)$ で割ると
\begin{equation}
  \frac{\dot{x}(t)}{x(t)}=-a
\end{equation}
である。合成関数の微分公式より、$\log |x(t)|$ の微分は
\begin{equation}
  \frac{\dd}{\dd t}\log |x(t)|=\frac{\dot{x}(t)}{x(t)}
\end{equation}
である。したがって
\begin{equation}
  \frac{\dd}{\dd t}\log |x(t)|=-a
\end{equation}
となる。時刻 $0$ から $t$ まで積分して
\begin{align}
  \int_0^t\frac{\dd}{\dd \tau}\log |x(\tau)|\,\dd\tau
  &=\int_0^t -a\,\dd\tau,\\
  \log |x(t)|-\log |x(0)|&=-at.
\end{align}
対数の性質 $\log A-\log B=\log(A/B)$ より
\begin{equation}
  \log \left|\frac{x(t)}{x(0)}\right|=-at
\end{equation}
である。両辺の指数をとれば
\begin{equation}
  \left|\frac{x(t)}{x(0)}\right|=e^{-at}
\end{equation}
となる。初期値の符号は解の途中で変わらないので、符号を含めて
\begin{equation}
  x(t)=x(0)e^{-at}
\end{equation}
を得る。
\end{derivation}

この導出を見ると、指数関数は経験的に仮定する形ではない。「変化率が現在値に比例する」という仮定から、対数の微分を通って出てくる必然的な形である。

\subsubsection{一次線形微分方程式の初期値問題}

制御対象の最小モデルは、減衰だけを表す \weq{basic-decay-ode} から、入力や外乱を受ける一次線形微分方程式へ広がる。係数 $a$ [1/s] を一定、右辺の信号 $r(t)$ の単位を $x$/s として
\begin{equation}
  \dot{x}(t)+a x(t)=r(t)
  \label{eq:first-order-linear-ivp}
\end{equation}
を考える。微分方程式だけでは解は一つに決まらない。時刻 $0$ の値
\begin{equation}
  x(0)=x_0
  \label{eq:first-order-linear-initial}
\end{equation}
を合わせて指定したものを初期値問題という。初期値は、コンデンサに最初から蓄えられている電荷、物体の初期温度、台車の初期位置や速度のように、入力を加える前から対象の内部に残っている情報を表す。

\begin{theorem}[一次線形微分方程式の積分因子解]
$a$ を定数、$r(t)$ を区分的に連続な信号とする。\weq{first-order-linear-ivp} と \weq{first-order-linear-initial} の解は
\begin{equation}
  x(t)=e^{-at}x_0+\int_0^t e^{-a(t-\tau)}r(\tau)\,\dd\tau
  \label{eq:first-order-integrating-factor-solution}
\end{equation}
である。
\end{theorem}

\begin{derivation}
\weq{first-order-linear-ivp} の両辺に $e^{at}$ を掛ける。
\begin{equation}
  e^{at}\dot{x}(t)+a e^{at}x(t)=e^{at}r(t)
\end{equation}
左辺は積の微分公式により
\begin{equation}
  \frac{\dd}{\dd t}\{e^{at}x(t)\}=e^{at}\dot{x}(t)+a e^{at}x(t)
\end{equation}
である。したがって
\begin{equation}
  \frac{\dd}{\dd t}\{e^{at}x(t)\}=e^{at}r(t)
\end{equation}
を得る。ここで $e^{at}$ を積分因子という。時刻 $0$ から $t$ まで積分すると
\begin{align}
  e^{at}x(t)-x(0)&=\int_0^t e^{a\tau}r(\tau)\,\dd\tau,\\
  x(t)&=e^{-at}x_0+\int_0^t e^{-a(t-\tau)}r(\tau)\,\dd\tau
\end{align}
となる。
\end{derivation}

\weq{first-order-integrating-factor-solution} の第一項は初期値から生じる応答であり、入力を 0 とした同次方程式
\begin{equation}
  \dot{x}+a x=0
\end{equation}
の解である。この応答を自然応答または自由応答という。特性方程式は $s+a=0$ であり、根 $s=-a$ に対応する固有モードは $e^{-at}$ である。したがって一次系の固有応答は指数関数一つで決まる。第二項は右辺信号 $r(t)$ によって作られる応答であり、強制応答という。線形方程式では、全応答は
\begin{equation}
  x(t)=x_{\mathrm{free}}(t)+x_{\mathrm{forced}}(t)
\end{equation}
のように、同次解と特殊解の和として表される。

右辺が一定値 $r(t)=r_0$ の場合は、特殊解を定数 $x_p$ として求められる。$\dot{x}_p=0$ を \weq{first-order-linear-ivp} に代入すると
\begin{equation}
  a x_p=r_0,\qquad x_p=\frac{r_0}{a}
  \label{eq:first-order-particular-constant}
\end{equation}
である。ただし $a\ne0$ とする。偏差 $z=x-x_p$ を置くと
\begin{equation}
  \dot{z}=\dot{x},\qquad
  \dot{z}+a z=0
\end{equation}
となる。変数分離法により
\begin{equation}
  \frac{\dot{z}}{z}=-a
\end{equation}
を積分して
\begin{equation}
  z(t)=z(0)e^{-at}
\end{equation}
を得る。したがって
\begin{equation}
  x(t)=\frac{r_0}{a}+
  \left(x_0-\frac{r_0}{a}\right)e^{-at}
  \label{eq:first-order-constant-input-solution}
\end{equation}
である。この式では $r_0/a$ が特殊解、括弧内の係数をもつ指数項が同次解である。

\begin{definition}[時定数]
$a>0$ の一次安定系では
\begin{equation}
  \tau=\frac{1}{a}
\end{equation}
を時定数という。時定数の単位は s であり、偏差 $z(t)$ は
\begin{equation}
  z(\tau)=z(0)e^{-1}\simeq0.368z(0)
\end{equation}
まで小さくなる。
\end{definition}

時定数は「何秒で完全に定常値へ到達するか」ではなく、偏差が同じ比率で減る速さを表す量である。$t=2\tau$ では $e^{-2}\simeq0.135$、$t=3\tau$ では $e^{-3}\simeq0.050$ になる。制御設計で応答を速くするとは、多くの場合、閉ループ方程式の係数を変えて固有モードの時定数を小さくすることを意味する。

\begin{example}[比例制御された一次対象の初期値応答]
一次対象
\begin{equation}
  \tau_p\dot{y}(t)+y(t)=K_0u(t),\qquad \tau_p>0
\end{equation}
を、比例制御器
\begin{equation}
  u(t)=k_c\{r(t)-y(t)\}
\end{equation}
で制御する。ここで $y$ は出力、$u$ は入力、$r$ は参照信号、$K_0$ は対象の静的ゲイン、$k_c$ は比例ゲインである。参照信号を一定値 $r(t)=r_0$ とすると
\begin{align}
  \tau_p\dot{y}+y&=K_0k_c(r_0-y),\\
  \tau_p\dot{y}+(1+K_0k_c)y&=K_0k_c r_0.
\end{align}
$1+K_0k_c>0$ のとき閉ループは一次安定系であり、
\begin{equation}
  a_{\mathrm{cl}}=\frac{1+K_0k_c}{\tau_p},
  \qquad
  r_{\mathrm{cl}}=\frac{K_0k_c}{\tau_p}r_0
\end{equation}
と置けば $\dot{y}+a_{\mathrm{cl}}y=r_{\mathrm{cl}}$ である。定常値と時定数は
\begin{equation}
  y_\infty=\frac{K_0k_c}{1+K_0k_c}r_0,
  \qquad
  \tau_{\mathrm{cl}}=\frac{1}{a_{\mathrm{cl}}}
  =\frac{\tau_p}{1+K_0k_c}
\end{equation}
となる。初期値を $y(0)=y_0$ とすれば
\begin{equation}
  y(t)=y_\infty+(y_0-y_\infty)e^{-t/\tau_{\mathrm{cl}}}
  \label{eq:proportional-first-order-response}
\end{equation}
である。\weq{proportional-first-order-response} の指数項は閉ループの自然応答であり、$y_\infty$ は一定参照信号に対する強制応答である。比例ゲイン $k_c$ を大きくすると $\tau_{\mathrm{cl}}$ は小さくなるが、定常値は一般に $r_0$ と一致しない。この残りの偏差を消すには、後で扱う積分動作やフィードフォワードの設計が必要になる。
\end{example}

\subsubsection{二次微分方程式への接続}

機械系では、一つの状態だけでは未来を決められないことが多い。質量ばねダンパ系の変位 $q(t)$ は、入力力 $u(t)$ と外力外乱を無視すれば
\begin{equation}
  M\ddot{q}(t)+c\dot{q}(t)+kq(t)=u(t)
  \label{eq:second-order-ode-connection}
\end{equation}
のような二次微分方程式に従う。$M$ [kg]、$c$ [N\,s/m]、$k$ [N/m] は正のパラメータである。一定入力 $u(t)=u_0$ に対する特殊解は
\begin{equation}
  q_p=\frac{u_0}{k}
\end{equation}
であり、偏差 $z=q-q_p$ は
\begin{equation}
  M\ddot{z}+c\dot{z}+kz=0
\end{equation}
を満たす。指数関数形 $z=e^{st}$ を代入すると
\begin{equation}
  Ms^2+cs+k=0
  \label{eq:mass-spring-characteristic}
\end{equation}
を得る。\weq{mass-spring-characteristic} が二次系の特性方程式であり、その二つの根が固有モードを決める。根が相異なる実数 $s_1,s_2$ なら固有応答は
\begin{equation}
  z(t)=C_1e^{s_1t}+C_2e^{s_2t}
\end{equation}
であり、複素共役根なら減衰振動として表される。係数 $C_1,C_2$ は、初期変位 $q(0)$ と初期速度 $\dot{q}(0)$ の二つの初期値から決まる。

この二次方程式も、状態を
\begin{equation}
  x(t)=\begin{bmatrix}q(t)\\ \dot{q}(t)\end{bmatrix}
\end{equation}
と選べば一階の連立微分方程式になる。つまり一次系で見た同次解、特殊解、初期値、時定数、固有応答という考えは、二次系では「固有根が一つから二つへ増える」形で拡張される。後の章で扱う行列指数関数は、この連立一次系の自然応答をまとめて書くための記法である。

\subsubsection{テイラー展開による局所近似}

制御器が利用できる情報は、現在値、過去の測定値、既知の入力に限られる。それでも、現在の値と変化率が分かれば、短い時間後の値は近似できる。

\begin{theorem}[テイラーの定理の一次近似]
十分なめらかな関数 $x(t)$ について、$\Delta t$ が小さいとき
\begin{equation}
  x(t+\Delta t)=x(t)+\dot{x}(t)\Delta t+O(\Delta t^2)
\end{equation}
である。二次まで残すと
\begin{equation}
  x(t+\Delta t)=x(t)+\dot{x}(t)\Delta t+
  \frac{1}{2}\ddot{x}(t)\Delta t^2+O(\Delta t^3)
\end{equation}
である。
\end{theorem}

\begin{derivation}
テイラーの定理より、$t$ のまわりの展開は
\begin{equation}
  x(t+\Delta t)=x(t)+x'(t)\Delta t+
  \frac{x''(t)}{2!}\Delta t^2+
  \frac{x^{(3)}(t)}{3!}\Delta t^3+\cdots
\end{equation}
である。ここで $x'(t)=\dot{x}(t)$、$x''(t)=\ddot{x}(t)$ と書き直す。一次近似では $\Delta t^2$ 以上の項をまとめて $O(\Delta t^2)$ と書くので、最初の式を得る。二次近似では $\Delta t^3$ 以上の項をまとめるので、二つ目の式を得る。
\end{derivation}

この定理は、線形化、数値積分、離散化の土台になる。小さい $\Delta t$ では $\Delta t^2$ や $\Delta t^3$ はさらに小さくなるため、短い時間だけなら最初の数項で動きを近似できる。

\subsubsection{陽的 Euler 法による数値積分と刻み幅}

連続時間モデル
\begin{equation}
  \dot{x}(t)=f(x(t),u(t),d(t),p,t)
  \label{eq:continuous-state-model-for-integration}
\end{equation}
は、各時刻における状態の変化率を与える。計算機シミュレーションでは、連続なすべての時刻をそのまま扱わず、時刻列
\begin{equation}
  t_k=t_0+kh,\qquad k=0,1,2,\ldots
\end{equation}
上の近似値 $x_k\simeq x(t_k)$ を順に更新する。ここで $h>0$ は刻み幅またはシミュレーションステップサイズであり、単位は s である。$x_i$ の単位が温度 K、電圧 V、位置 m のいずれであっても、$f_i$ の単位は $x_i$/s でなければならない。したがって $h f_i$ は $x_i$ と同じ単位をもち、状態更新で足し合わせられる量になる。

\begin{formula}[陽的 Euler 法]
\weq{continuous-state-model-for-integration} において、区間 $[t_k,t_{k+1})$ で入力と外乱を代表値 $u_k$、$d_k$ により近似する。$f$ が十分なめらかで、刻み幅 $h$ が小さいとき、陽的 Euler 法は
\begin{equation}
  x_{k+1}=x_k+h f(x_k,u_k,d_k,p,t_k)
  \label{eq:explicit-euler-state-update}
\end{equation}
により状態を更新する。
\end{formula}

\begin{derivation}
微分方程式を時刻 $t_k$ から $t_{k+1}=t_k+h$ まで積分すると、厳密解は
\begin{equation}
  x(t_{k+1})=x(t_k)+
  \int_{t_k}^{t_k+h} f(x(\tau),u(\tau),d(\tau),p,\tau)\,\dd\tau
  \label{eq:exact-integral-state-update}
\end{equation}
を満たす。陽的 Euler 法では、この積分区間の被積分関数を左端の値
\begin{equation}
  f(x(t_k),u_k,d_k,p,t_k)
\end{equation}
で一定と近似する。すると
\begin{align}
  x(t_{k+1})
  &\simeq x(t_k)+
  \int_{t_k}^{t_k+h} f(x(t_k),u_k,d_k,p,t_k)\,\dd\tau\\
  &=x(t_k)+h f(x(t_k),u_k,d_k,p,t_k)
\end{align}
となる。厳密値 $x(t_k)$ の代わりに近似値 $x_k$ を使うと \weq{explicit-euler-state-update} を得る。
\end{derivation}

この更新式は、微分を差分で置き換えただけの記法ではなく、現在の変化率を時間 $h$ だけ積み上げる近似である。一回の更新を厳密値 $x(t_k)$ から開始した場合、テイラー展開により
\begin{align}
  x(t_k+h)
  &=x(t_k)+h\dot{x}(t_k)+\frac{h^2}{2}\ddot{x}(\xi_k)\\
  &=x(t_k)+h f(x(t_k),u_k,d_k,p,t_k)+O(h^2)
\end{align}
と書ける。ただし $\xi_k$ は $t_k$ と $t_k+h$ の間の時刻であり、区間内で入力と外乱を一定値で近似しているとする。したがって一ステップの局所離散化誤差は $O(h^2)$ である。一方、固定した時間区間 $[t_0,T]$ を $N=(T-t_0)/h$ ステップで進むと、局所誤差が各ステップで蓄積する。$f$ が状態に関して Lipschitz 連続で、必要な導関数が有界であるなら、陽的 Euler 法の大域誤差は $O(h)$ で抑えられる。刻み幅を半分にすると、同じ時間区間のステップ数はほぼ二倍になり、大域誤差は一次の精度で小さくなる。

刻み幅は精度だけでなく数値安定性にも関わる。最も単純な減衰方程式
\begin{equation}
  \dot{z}(t)=-a z(t),\qquad a>0
\end{equation}
に陽的 Euler 法を適用すると
\begin{equation}
  z_{k+1}=(1-ah)z_k
  \label{eq:euler-stability-test}
\end{equation}
である。連続時間では任意の $a>0$ で $z(t)$ は 0 へ減衰する。しかし離散更新では
\begin{equation}
  |1-ah|<1
\end{equation}
すなわち
\begin{equation}
  0<h<\frac{2}{a}
  \label{eq:euler-stability-step-bound}
\end{equation}
でなければ、数列 $z_k$ は減衰しない。$h$ が大きすぎると、連続時間モデルは安定であるにもかかわらず、数値解だけが符号反転を伴って増大することがある。これは対象の物理的な不安定性ではなく、数値積分法と刻み幅の組合せから生じる数値不安定性である。

RC 回路
\begin{equation}
  C_e\dot{v}_C(t)=\frac{u(t)-v_C(t)}{R}
\end{equation}
で、抵抗 $R$ [$\Omega$]、容量 $C_e$ [F] を一定、入力電圧を区間ごとに $u_k$ [V] で一定と仮定する。時定数を
\begin{equation}
  \tau=R C_e
\end{equation}
と置くと、状態 $x=v_C$ は
\begin{equation}
  \dot{x}(t)=-\frac{1}{\tau}x(t)+\frac{1}{\tau}u(t)
\end{equation}
に従う。陽的 Euler 法では
\begin{align}
  x_{k+1}
  &=x_k+h\left(-\frac{1}{\tau}x_k+\frac{1}{\tau}u_k\right)\\
  &=\left(1-\frac{h}{\tau}\right)x_k+\frac{h}{\tau}u_k
  \label{eq:rc-explicit-euler-update}
\end{align}
である。比 $h/\tau$ は無次元であり、電圧 $x_k$ と $u_k$ に掛けられる係数として解釈できる。入力を 0 とした自然応答では、数値的な減衰条件は
\begin{equation}
  0<h<2\tau
\end{equation}
である。たとえば $\tau=0.5\,\mathrm{s}$ なら $h=0.01\,\mathrm{s}$ は時定数に比べて十分小さいが、$h=1.5\,\mathrm{s}$ は \weq{euler-stability-step-bound} を破り、連続時間の RC 回路とは異なる発散的な数列を作る。後の状態空間離散化や実装章で扱うサンプリング周期は、この刻み幅を制御器、センサ、アクチュエータ、数値積分法の条件と合わせて選ぶ問題である。
